
\begingroup
\shorthandoff{<>}% Babel español trata de otro modo ciertos delimitadores de Xy-pic.
\setcounter{subsection}{9}% Conservar el número visible XIII.2.10 de la fuente.
\subsection{}
\label{XIII.2.10}
Si $U$ es un esquema conexo, $a$ un punto geométrico de~$U$ y
$\LL$ un conjunto de números primos, se denota por
\begin{equation*}
\label{eq:XIII.2.10.0}
\tag{2.10.0} \pi_1^{\LL}(U,a)
\end{equation*}
\label{indnot:mi}\oldindexnot{$\pi_1^{\LL}(U,a)$|hyperpage}%
el límite proyectivo de los cocientes finitos de~$\pi_1(U,a)$ cuyos
órdenes tienen todos sus factores primos en~$\LL$.

Vamos a definir homomorfismos de especialización para el grupo
fundamental que generalizan X.\Ref{X.2}.

Sea $g\colon U\to S$ un morfismo coherente de fibras geométricamente
conexas (\resp un morfismo de la forma $g=fi$, donde $f\colon X\to S$
es un morfismo propio de presentación finita e $i\colon U\to X$ es una
inmersión abierta tal que $U$ sea el complementario en~$X$ de un divisor
de cruzamientos normales relativamente a~$S$ (\cf \Ref{XIII.2.8})). Sea
$\LL$ un conjunto de números primos y supongamos, en el caso sin
paréntesis, que, para todo $\LL$-grupo constante finito $C$, $(C_U,g)$
sea cohomológicamente propio relativamente a~$S$ en dimensión $\leq 1$.
Sea $\overline{s}_1\to\overline{s}_2$ un morfismo de especialización de
puntos geométricos de~$S$, sea $\overline{S}$ la localización estricta
de~$S$ en $\overline{s}_2$ y sea
$\overline{U}=U\times_S\overline{S}$. Se tiene un diagrama conmutativo
$$
\ifthenelse{\boolean{orig}}
{\xymatrix{ U_{\bar{s}_1} \ar[r]^{h_1} \ar[d] & \bar{U}
\ar[d]^{\bar{g}} & U_{\bar{s}_1} \ar[l]_{h_2} \ar[d] \\ \bar{s}_1
\ar[r] & \bar{S} & \,\bar{s}_2\,. \ar[l] }}
{\xymatrix{ U_{\bar{s}_1} \ar[r]^{h_1} \ar[d] & \bar{U}
\ar[d]^{\bar{g}} & U_{\bar{s}_2} \ar[l]_{h_2} \ar[d] \\ \bar{s}_1
\ar[r] & \bar{S} & \,\bar{s}_2\,. \ar[l] }}
$$
Si $a_1$ es un punto geométrico de~$U_{\bar{s}_1}$ y $a_2$ un
punto geométrico de~$U_{\bar{s}_2}$, los morfismos $h_1$ y $h_2$
definen homomorfismos canónicos
\marginpar{391}%
\begin{align*}
\pi_1\colon \pi_1^{\LL}(U_{\bar{s}_1},a_1)&\to
\pi_1^{\LL}(\bar{U},a_1) & \pi_2\colon
\pi_1^{\LL}(U_{\bar{s}_2},a_2)&\to \pi_1^{\LL}(\bar{U},a_2)\\
(\text{\resp} \pi_1\colon \pi_1^\tame (U_{\bar{s}_1},a_1)&\to
\pi_1^\tame (\bar{U},a_1) & \pi_2\colon
\pi_1^\tame(U_{\bar{s}_2},a_2)&\to \pi_1^\tame(\bar{U},a_2))
\end{align*}
(\Ref{V.7} y~\Ref{eq:XIII.2.1.5.2}). Las hipótesis relativas a la
propiedad de ser cohomológicamente propio (\resp
\ifthenelse{\boolean{orig}}
{\Ref{XIII.2.8}}
{\Ref{XIII.2.7}})
demuestran que $\pi_2$ es un isomorfismo. Si se elige una clase de
caminos de~$a_1$ a~$a_2$,
se obtiene un isomorfismo
$$
\pi_{12}\colon \pi_1^{\LL}(\bar{U},a_1)\isomto
\pi_1^{\LL}(\bar{U},a_2) \quad (\text{\resp} \pi_{12}\colon
\pi_1^\tame (\bar{U},a_1)\isomto \pi_1^\tame(\bar{U},a_2))
$$
de donde resulta un homomorfismo $\pi=\pi_2^{-1}\pi_{12}\pi_1$
$$
\pi\colon \pi_1^{\LL}(U_{\bar{s}_1},a_1)\to
\pi_1^{\LL}(U_{\bar{s}_2},a_2) \quad (\text{\resp}\pi\colon
\pi_1^\tame (U_{\bar{s}_1},a_1)\to \pi_1^\tame
(U_{\bar{s}_2},a_2)).
$$
Cambiar la clase de caminos de~$a_1$ a~$a_2$ equivale a modificar
$\pi$ por un automorfismo interior de
\ifthenelse{\boolean{orig}}
{$\pi_1^{\LL}(X_{\bar{s}_2},a_2)$ (\resp de
$\pi_1^\tame(X_{\bar{s}_2},a_2)$)}
{$\pi_1^{\LL}(U_{\bar{s}_2},a_2)$ (\resp de
$\pi_1^\tame(U_{\bar{s}_2},a_2)$)}. Se llama
\emph{homomorfismo de especialización para el grupo fundamental}
\index{homomorfismo de especialización para el grupo fundamental|hyperpage}%
\index{especialización para el grupo fundamental (homomorfismo de)|hyperpage}%
asociado al morfismo $\bar{s}_1\to\bar{s}_2$ y se denota simplemente por
$$
\ifthenelse{\boolean{orig}}
{\pi\colon \pi_1^{\LL}(X_{\bar{s}_1})\to
\pi_1^{\LL}(X_{\bar{s}_2}) \quad \text{(\resp} \pi\colon
\pi_1^{\tame}(X_{\bar{s}_1})\to \pi_1^{\tame}(X_{\bar{s}_2}))}
{\pi\colon \pi_1^{\LL}(U_{\bar{s}_1})\to
\pi_1^{\LL}(U_{\bar{s}_2}) \quad \text{(\resp} \pi\colon
\pi_1^{\tame}(U_{\bar{s}_1})\to \pi_1^{\tame}(U_{\bar{s}_2}))}
$$
uno de los homomorfismos definidos anteriormente.

\endgroup
